
\subsubsection{Deflection and Rotation Fields}

We introduce three classes of gauge parameters which represent the ambiguity in whether terms in \(\mathcal{R}(\Section)\) are assigned to the beam fields \(\bm{u},\bm{\theta}\) or the warping \(\bm{\alpha}\):

\begin{enumerate}
\item \textbf{Shear gauge} $\bm{\eta}_{\Section} \in \mathbb{R}^2$: Shifts the distribution of transverse rotation between $\bm{\theta}$ and compensating warping modes
\item \textbf{Axial gauge} $\eta_\varepsilon \in \mathbb{R}$: Shifts constant axial strain between $\bm{u}$ and axial warping
\item \textbf{Twist gauge} $\eta_{\vartheta} \in \mathbb{R}$: Redistributes twist between $\bm{\theta}$ and torsional warping
\item \textbf{?gauge} $\bm{\eta}_{v}$
\end{enumerate}

The gauged kinematic fields become:
%
\begin{equation}
\begin{aligned}
\bm{u}(x) &= x(a_\varepsilon + \eta_\varepsilon)\FrameBasis
+ \reflenx\,\GaugeShearB
- \tfrac{1}{2}x^2\IesanLoadM - \tfrac{1}{6}x^3\bm{b}_{\Section} - \tfrac{1}{2}x^2(\FrameBasis\otimes\CentroidLocation)\bm{b}_{\Section} \\
\bm{\theta}(x) &= \GaugeTwist(x)\FrameBasis - \FrameBasis \times \GaugeShearA + x(a_\vartheta + c_\vartheta)\FrameBasis - x(\FrameBasis \times \IesanLoadM) - \tfrac{1}{2}x^2(\FrameBasis \times \bm{b}_{\Section})
\end{aligned}
\label{eq:gauged-kinematics}
\end{equation}

Each gauge parameter requires a compensating warping mode to maintain exactness with the Saint-Venant solution:
\begin{subequations}
\begin{align*}
\bm{\Psi}_\varepsilon(\bm{r}) &= \FrameBasis,  &\alpha_\varepsilon &= -\reflenx\,\eta_\varepsilon \\
\bm{\Psi}_{\Section}(\bm{r}) &= \oneS &\bm{\alpha}_{.} &= -x\GaugeShearB \\
\bm{\Psi}_\vartheta(\bm{r}) &= \FrameBasis \times \bm{r},  &\GaugeTwist(x) &= -\reflenx\,\bar{\GaugeTwist} \\
\bm{\Psi}_r(\bm{r}) &= \FrameBasis \otimes \bm{r},
&\bm{\alpha}_r &= -\GaugeShearA \\
\end{align*}
\end{subequations}
These compensating modes exactly cancel the spurious displacements introduced by the gauge shift. 
This choice is what distinguishes the  ``energetic'' identification of \cite{renton1991generalized} from
 the geometric identification of \cite{cowper1966shear}. 
 The total displacement
\(\hat{\bm{u}}\) remains exact regardless of this choice -- only the
decomposition into beam kinematics changes.
%

From \eqref{eq:sv-kinematics-ii}, 1D deformations \(\bm{\gamma}=\bm{u}'+\FrameBasis\times\bm{\theta}\) and \(\bm{\kappa}=\bm{\theta}'\) obtained from \eqref{eq:gauged-kinematics} are:

\begin{equation}
\begin{aligned}
\bm{\gamma}(x) 
% &= \bm{u}' + \FrameBasis \times \bm{\theta} \\
% &= (a_\varepsilon + \eta_\varepsilon)\FrameBasis 
%  - \reflenx\IesanLoadM - \tfrac{1}{2}x^2\bm{b}_{\Section} - x(\FrameBasis\otimes\CentroidLocation)\bm{b}_{\Section} \\
% &\quad + \FrameBasis \times (\eta_\vartheta\FrameBasis - \FrameBasis \times \GaugeShearA) + x\IesanLoadM + \tfrac{1}{2}x^2\bm{b}_{\Section} \\
&= (a_\varepsilon + \eta_\varepsilon - x\CentroidLocation\cdot\bm{b}_{\Section})\FrameBasis
+ \bm{\eta}_{\Section} + \bm{\eta}_v \\
%
\bm{\kappa}(x) &= (a_\vartheta + c_\vartheta + \GaugeTwist')\FrameBasis - (\FrameBasis \times \IesanLoadM) - x(\FrameBasis \times \bm{b}_{\Section})
\end{aligned}
\label{eq:gauged-deformations}
\end{equation}

The gauge parameter $\bm{\eta}_{\Section}$ now appears directly in the strain measure, providing control over the shear strain distribution. 
